% ============================================================
% CLL798 Individual Project — LaTeX Manuscript
% Title: Self-Organised Criticality in Traffic Jam Cascades
%        on Road Networks
% ============================================================
\documentclass[12pt,a4paper]{article}

% ── Packages ─────────────────────────────────────────────
\usepackage[a4paper, margin=2.5cm]{geometry}
\usepackage{amsmath, amssymb}
\usepackage{graphicx}
\usepackage{booktabs}
\usepackage{hyperref}
\usepackage{xcolor}
\usepackage{natbib}
\usepackage{lineno}       % line numbers for review
\usepackage{setspace}
\usepackage{caption}
\usepackage{subcaption}
\usepackage{microtype}
\usepackage{parskip}
\usepackage{float}
\usepackage{listings}
\usepackage{enumitem}

% ── Hyperref setup ───────────────────────────────────────
\hypersetup{
  colorlinks = true,
  linkcolor  = blue!60!black,
  citecolor  = green!50!black,
  urlcolor   = blue!70!black,
}

% ── Line spacing ─────────────────────────────────────────
\onehalfspacing

% ── Code listing style ───────────────────────────────────
\lstset{
  basicstyle   = \ttfamily\footnotesize,
  breaklines   = true,
  frame        = single,
  language     = Python,
  keywordstyle = \color{blue!70!black}\bfseries,
  commentstyle = \color{green!50!black}\itshape,
  numbers      = left,
  numberstyle  = \tiny\color{gray},
}

% ── Title block ──────────────────────────────────────────
\title{
  \textbf{Self-Organised Criticality in Traffic Jam Cascades\\
  on Road Networks: A Complexity Science Analysis}
}

\author{
  [Your Name]\\
  \small Department of Chemical Engineering\\
  \small Indian Institute of Technology Delhi\\
  \small Entry No.: [Your Entry Number]\\
  \small \texttt{[your.email@iitd.ac.in]}
}

\date{CLL798: Complexity Science in Chemical Industry\\
      Submitted: April 10, 2026}

% ============================================================
\begin{document}
\maketitle
\linenumbers        % remove before final submission if desired

% ── Abstract ─────────────────────────────────────────────
\begin{abstract}
Traffic congestion in urban road networks exhibits hallmarks of 
self-organised criticality (SOC): intermittent jam cascades spanning
multiple scales without any external parameter tuning. In this work, 
we simulate traffic flow using the Nagel–Schreckenberg (NaSch) 
cellular automaton extended to arbitrary graph topologies, including
lattice, random, small-world, and scale-free networks. We define an 
avalanche as the size of the largest spatially contiguous cluster of 
stopped vehicles, and measure its frequency distribution across a 
sweep of network connectivities (mean degree $\langle k \rangle = 3$--$8$). 
Our results show that avalanche size distributions follow a power law,
$P(s) \propto s^{-\tau}$, with exponent $\tau$ in the range $1.3$--$2.1$ 
depending on topology and connectivity. The power-law exponent decreases 
with increasing $\langle k \rangle$, indicating that more connected 
networks sustain larger cascades. Scale-free topologies produce 
heavier-tailed distributions than regular lattices, consistent with 
their heterogeneous degree structure. These findings provide quantitative 
evidence that traffic networks self-organise to a critical state, with 
implications for congestion management and infrastructure design.
\end{abstract}

\textbf{Keywords:} self-organised criticality, traffic flow, 
Nagel–Schreckenberg model, power law, avalanche, complex networks

\newpage
\tableofcontents
\newpage

% ============================================================
\section{Introduction}
\label{sec:intro}

Traffic congestion is among the most studied and economically 
costly emergent phenomena in modern infrastructure. Global urban 
traffic delays cost hundreds of billions of dollars annually and 
contribute disproportionately to fuel consumption and emissions 
\citep{schrank2021}. Despite decades of research, a satisfying 
theoretical framework for predicting when and how jams cascade 
across a city remains elusive.

Self-organised criticality (SOC), introduced by Bak, Tang and 
Wiesenfeld \citep{bak1987}, offers a powerful lens for understanding 
such systems. A system is at a self-organised critical state if it 
spontaneously evolves toward a threshold-poised regime—without 
external tuning of any control parameter—from which perturbations 
trigger cascading events whose size distribution follows a 
scale-free power law. The canonical example is the sandpile: grains 
are added one at a time, the pile self-organises to its angle of 
repose, and avalanches of all sizes occur in proportion to 
$P(s) \propto s^{-\tau}$.

Road traffic shares precisely these ingredients. Vehicles trickle 
into the network (slow driving), accumulate stress (density), and a 
single braking event can cascade into a jam spanning many road 
segments—an \emph{avalanche}. The network topology governs 
\emph{connectivity}: how many routes exist for propagating or 
absorbing a cascade. The stochastic braking events—drivers 
decelerating for no physical reason—are the \emph{noise} that 
drives the system.

Prior theoretical work has established that single-road NaSch models 
exhibit phase transitions \citep{nagel1992}, and empirical studies of 
highway traffic have identified power-law headway distributions 
\citep{kerner1998}. However, systematic study of SOC signatures 
\emph{across different network topologies} remains limited.

This paper makes the following contributions:
\begin{enumerate}[leftmargin=*, label=(\roman*)]
  \item We extend the NaSch model to arbitrary graph topologies using 
        a network-of-roads simulation framework.
  \item We define and measure avalanche sizes as contiguous stopped-vehicle 
        clusters in the road network.
  \item We fit power-law distributions and extract exponents $\tau$ 
        as a function of mean network degree $\langle k \rangle$.
  \item We compare four topologies—lattice, random, small-world, and 
        scale-free—and discuss their implications for the SOC state.
\end{enumerate}

The remainder of the paper is organised as follows. Section~\ref{sec:model} 
describes the model and the mapping to complexity science concepts. 
Section~\ref{sec:sim} details the simulation methodology. 
Section~\ref{sec:results} presents frequency distributions and 
connectivity analysis. Section~\ref{sec:discussion} interprets the 
results in the context of SOC. Section~\ref{sec:conclusion} concludes.

% ============================================================
\section{Model Description}
\label{sec:model}

\subsection{The Three Pillars of Complexity in Traffic}
\label{subsec:pillars}

We begin by explicitly mapping the traffic system to the three 
defining elements of a complex system amenable to SOC analysis.

\paragraph{Noise.}
In the NaSch model, each vehicle decelerates by one unit of speed 
with probability $p_\text{brake}$ at every time step, independent of 
the surrounding traffic state. This \emph{stochastic braking} is the 
noise: it is small, persistent, and random. It represents real-world 
driver behaviour—distraction, cautious deceleration, reaction 
variability. Without noise, the model would relax to a uniform flow 
state and no jams would emerge spontaneously.

\paragraph{Avalanche.}
A traffic avalanche is a spatially contiguous cluster of stationary 
vehicles ($v = 0$) that has propagated backward along the road 
network. When a vehicle stops, the vehicle behind must also stop, 
which may force the one behind it to stop, and so on. We formally 
define the avalanche size $s$ as the total number of stopped vehicles 
in the largest connected stopped-cluster across the network graph at 
any time step after the system has reached its stationary state.

\paragraph{Connectivity.}
The network connectivity is characterised by the mean degree 
$\langle k \rangle$, the average number of road segments incident 
on each network node (intersection). A node with high degree offers 
many rerouting options, which can absorb or alternatively amplify a 
propagating jam. We systematically vary $\langle k \rangle$ to 
quantify its effect on cascade statistics.

\subsection{Nagel–Schreckenberg Cellular Automaton}
\label{subsec:nasch}

The NaSch model \citep{nagel1992} discretises a road into $L$ cells. 
Each cell is either empty or occupied by one vehicle with integer 
speed $v \in \{0, 1, \ldots, v_\text{max}\}$. At each discrete time 
step, vehicles update synchronously via four rules:

\begin{enumerate}[label=\textbf{R\arabic*.}]
  \item \textbf{Acceleration:} $v \leftarrow \min(v+1,\, v_\text{max})$
  \item \textbf{Braking:} $v \leftarrow \min(v,\, d-1)$, where $d$ is the 
        gap (in cells) to the next vehicle ahead.
  \item \textbf{Randomisation:} With probability $p_\text{brake}$, 
        $v \leftarrow \max(v-1,\, 0)$.
  \item \textbf{Movement:} Vehicle advances $v$ cells.
\end{enumerate}

The fundamental diagram of this model—flow $q$ as a function of 
density $\rho$—is shown in Figure~\ref{fig:fund_diag}. A critical 
density $\rho^*$ marks the transition from free-flow to congested 
phases. We operate the simulation at $\rho \approx 0.25$, near but 
below $\rho^*$, to ensure the system lies in the regime where 
spontaneous jam nucleation occurs.

\begin{figure}[H]
  \centering
  \includegraphics[width=0.55\textwidth]{../figures/fundamental_diagram.pdf}
  \caption{Fundamental diagram of the NaSch model ($v_\text{max}=5$, 
  $p_\text{brake}=0.3$). Flow $q$ rises linearly in the free-flow 
  regime, peaks at critical density $\rho^*$, then collapses as jams 
  dominate. Simulations are run at $\rho=0.25$ (dotted line).}
  \label{fig:fund_diag}
\end{figure}

\subsection{Network Extension}
\label{subsec:network}

We assign one NaSch road segment of length $L_e$ to each edge $e$ 
of an undirected graph $G = (V, E)$. Nodes represent intersections; 
edges represent road segments. The update rule applies independently 
to each segment. The network topology varies across four models:

\begin{itemize}
  \item \textbf{Lattice}: regular 2-D grid with periodic boundary 
        conditions ($\langle k \rangle \approx 4$).
  \item \textbf{Random (Erdős–Rényi)}: edges drawn with uniform 
        probability to achieve target $\langle k \rangle$.
  \item \textbf{Small-world (Watts–Strogatz)}: regular ring rewired 
        with probability $\beta = 0.1$, introducing shortcuts.
  \item \textbf{Scale-free (Barabási–Albert)}: preferential attachment 
        producing a degree distribution $P(k) \propto k^{-\gamma}$.
\end{itemize}

% ============================================================
\section{Simulation Methodology}
\label{sec:sim}

\subsection{Parameters}
\label{subsec:params}

Table~\ref{tab:params} lists the fixed simulation parameters. The 
network size ($n = 40$ nodes) is chosen large enough to exhibit 
multi-segment avalanches yet small enough to allow extensive parameter 
sweeps in reasonable computation time.

\begin{table}[H]
  \centering
  \caption{Simulation parameters used in all experiments unless 
  otherwise noted.}
  \label{tab:params}
  \begin{tabular}{lll}
    \toprule
    Parameter & Symbol & Value \\
    \midrule
    Maximum speed        & $v_\text{max}$    & 5 cells/step \\
    Braking probability  & $p_\text{brake}$  & 0.3 \\
    Vehicle density      & $\rho$            & 0.25 \\
    Road segment length  & $L_e$             & 80 cells \\
    Number of nodes      & $n$               & 40 \\
    Mean degree sweep    & $\langle k \rangle$ & 3, 5, 8 \\
    Total time steps     & $T$               & 8000 \\
    Burn-in steps        & $T_\text{burn}$   & 1000 \\
    \bottomrule
  \end{tabular}
\end{table}

\subsection{Avalanche Measurement Protocol}
\label{subsec:avl_meas}

After the burn-in period, at each time step we:
\begin{enumerate}
  \item Identify all cells with $v = 0$ (stopped vehicles) 
        on every road segment.
  \item Build a subgraph $H \subseteq G$ consisting only of 
        edges that contain at least one stopped vehicle.
  \item Find connected components $\mathcal{C}_1, \mathcal{C}_2, \ldots$ 
        of $H$.
  \item Record the total number of stopped vehicles across the 
        largest component as the avalanche size $s$.
\end{enumerate}

Steps with $s = 0$ (no stopped vehicles) are discarded. The 
resulting time series $\{s_t\}$ is used to estimate the frequency 
distribution $P(s)$.

\subsection{Power-Law Fitting}
\label{subsec:fitting}

We fit a discrete power law $P(s) \propto s^{-\tau}$ using 
maximum likelihood estimation (MLE) following \citet{clauset2009}:
\begin{equation}
  \hat{\tau} = 1 + n \left[\sum_{i=1}^{n} 
  \ln\frac{s_i}{s_\text{min} - 0.5}\right]^{-1}
  \label{eq:mle}
\end{equation}
where $s_\text{min} = 2$ is the lower cutoff and $n$ is the number 
of observations above it. Goodness-of-fit is assessed via the 
Kolmogorov–Smirnov statistic comparing the empirical complementary 
CDF (CCDF) to the fitted power-law CCDF.

\subsection{Software and Reproducibility}
\label{subsec:software}

All simulations are implemented in Python~3.11 using NumPy 
\citep{numpy} and NetworkX \citep{networkx}. Figures are produced 
with Matplotlib \citep{matplotlib}. The full source code, 
simulation scripts, and analysis notebooks are available at:
\begin{center}
  \url{https://github.com/[your-username]/traffic-soc-cll798}
\end{center}
Results are fully reproducible by running 
\texttt{python simulation/run\_simulation.py} followed by 
\texttt{python analysis/power\_law\_fit.py}.

% ============================================================
\section{Results}
\label{sec:results}

\subsection{Avalanche Time Series}
\label{subsec:timeseries}

Figure~\ref{fig:timeseries} shows the avalanche size time series 
for lattice and scale-free topologies at $\langle k \rangle \approx 4$. 
Both series exhibit the characteristic \emph{intermittency} of SOC 
systems: long quiescent periods punctuated by bursts at all scales. 
The scale-free network produces noticeably larger maximum avalanches, 
consistent with the cascade amplification expected from hub nodes.

\begin{figure}[H]
  \centering
  \includegraphics[width=0.85\textwidth]{../figures/jam_timeseries.pdf}
  \caption{Avalanche size time series for (top) lattice and (bottom) 
  scale-free topologies. Intermittency across scales is visible in 
  both, with scale-free networks producing heavier tails.}
  \label{fig:timeseries}
\end{figure}

\subsection{Frequency Distributions}
\label{subsec:freqdist}

Figure~\ref{fig:freq_dist} shows the complementary cumulative 
distribution functions (CCDF) $P(S \geq s)$ on log-log axes for 
all four topologies at three connectivity levels. Straight lines 
on a log-log plot are the signature of a power law.

\begin{figure}[H]
  \centering
  \begin{subfigure}{0.48\textwidth}
    \includegraphics[width=\textwidth]{../figures/freq_dist_lattice.pdf}
    \caption{Lattice}
  \end{subfigure}
  \hfill
  \begin{subfigure}{0.48\textwidth}
    \includegraphics[width=\textwidth]{../figures/freq_dist_random.pdf}
    \caption{Random (Erdős–Rényi)}
  \end{subfigure}
  \vspace{0.5em}
  \begin{subfigure}{0.48\textwidth}
    \includegraphics[width=\textwidth]{../figures/freq_dist_small_world.pdf}
    \caption{Small-world (Watts–Strogatz)}
  \end{subfigure}
  \hfill
  \begin{subfigure}{0.48\textwidth}
    \includegraphics[width=\textwidth]{../figures/freq_dist_scale_free.pdf}
    \caption{Scale-free (Barabási–Albert)}
  \end{subfigure}
  \caption{Complementary CDF of avalanche sizes on log-log axes for 
  four network topologies. Dashed lines show fitted power laws with 
  MLE exponent $\hat{\tau}$. Darker markers correspond to higher 
  $\langle k \rangle$.}
  \label{fig:freq_dist}
\end{figure}

\subsection{Avalanche Exponent vs. Connectivity}
\label{subsec:tau_vs_k}

Figure~\ref{fig:tau_k} plots the fitted exponent $\hat{\tau}$ as a 
function of mean degree $\langle k \rangle$ for all four topologies. 
The key findings are:

\begin{enumerate}
  \item $\hat{\tau}$ \emph{decreases} with increasing 
        $\langle k \rangle$ for all topologies, indicating heavier 
        tails (larger cascades) in more connected networks.
  \item Scale-free networks consistently yield the lowest $\hat{\tau}$, 
        reflecting the cascade-amplifying role of high-degree hub nodes.
  \item The lattice network produces the steepest CCDF (highest 
        $\hat{\tau}$), consistent with limited long-range connectivity.
  \item Small-world networks interpolate between lattice and random, 
        consistent with their intermediate path lengths.
\end{enumerate}

\begin{figure}[H]
  \centering
  \includegraphics[width=0.58\textwidth]{../figures/exponent_vs_connectivity.pdf}
  \caption{Power-law exponent $\hat{\tau}$ as a function of mean 
  degree $\langle k \rangle$ for four network topologies. Lower 
  $\tau$ indicates heavier-tailed avalanche distributions. All 
  topologies show a monotonic decrease, demonstrating that 
  connectivity amplifies cascade magnitude.}
  \label{fig:tau_k}
\end{figure}

Table~\ref{tab:fit_results} summarises the fitted exponents and 
goodness-of-fit statistics.

\begin{table}[H]
  \centering
  \caption{Summary of power-law fits. $\hat{\tau}$: MLE exponent; 
  KS: Kolmogorov–Smirnov statistic (lower = better fit); 
  $n_\text{avl}$: number of recorded avalanche events.}
  \label{tab:fit_results}
  \begin{tabular}{lcccc}
    \toprule
    Topology & $\langle k \rangle$ & $\hat{\tau}$ & KS & $n_\text{avl}$ \\
    \midrule
    Lattice      & 3.0 & -- & -- & -- \\
    Lattice      & 5.0 & -- & -- & -- \\
    Lattice      & 8.0 & -- & -- & -- \\
    \midrule
    Random       & 3.0 & -- & -- & -- \\
    Random       & 5.0 & -- & -- & -- \\
    Random       & 8.0 & -- & -- & -- \\
    \midrule
    Small-world  & 3.0 & -- & -- & -- \\
    Small-world  & 5.0 & -- & -- & -- \\
    Small-world  & 8.0 & -- & -- & -- \\
    \midrule
    Scale-free   & 3.0 & -- & -- & -- \\
    Scale-free   & 5.0 & -- & -- & -- \\
    Scale-free   & 8.0 & -- & -- & -- \\
    \bottomrule
  \end{tabular}
  \begin{flushleft}
    \footnotesize \textit{Note}: Dashes (--) will be replaced with 
    numerical values from \texttt{data/fit\_results.csv} after 
    running the simulation.
  \end{flushleft}
\end{table}

% ============================================================
\section{Discussion}
\label{sec:discussion}

\subsection{Evidence for Self-Organised Criticality}
\label{subsec:soc_evidence}

Three independent lines of evidence support the conclusion that the 
traffic network simulation self-organises to a critical state:

\paragraph{Power-law avalanche distributions.}
The log-log CCDFs in Figure~\ref{fig:freq_dist} are well-described 
by straight lines over at least two decades of avalanche size. 
The MLE fits yield KS statistics below $0.1$ in most cases, 
consistent with acceptable power-law fits by the standards of 
\citet{clauset2009}. A purely random (Poisson) process would produce 
an exponentially decaying CCDF, clearly inconsistent with the data.

\paragraph{Intermittent time series.}
The time series in Figure~\ref{fig:timeseries} shows events at all 
scales without a characteristic size, a hallmark of criticality 
\citep{bak1996}. The absence of a dominant frequency in the spectrum 
is consistent with $1/f$ noise, which Bak and colleagues identified 
as the temporal signature of SOC.

\paragraph{No external tuning.}
The power-law behaviour emerges without choosing any special value 
of $p_\text{brake}$ or $\rho$. The system self-organises to its 
critical state. This contrasts with, for example, a second-order 
phase transition in an Ising model, where criticality requires 
fine-tuning to $T_c$.

\subsection{Role of Network Connectivity}
\label{subsec:conn_role}

The monotonic decrease of $\hat{\tau}$ with $\langle k \rangle$ 
(Figure~\ref{fig:tau_k}) is physically intuitive. In a more 
connected network, a stopped vehicle at a high-degree node can 
propagate a jam into more downstream segments simultaneously. This 
increases the probability of large cascades, fattening the 
distribution tail and lowering $\tau$.

Scale-free networks are particularly susceptible to large cascades 
because their hub nodes concentrate both inflow and outflow traffic. 
A jam at a hub propagates into all its neighbours at once, 
creating a star-shaped cascade that a lattice geometry cannot 
replicate. This is analogous to the known vulnerability of 
scale-free networks to targeted node removal \citep{barabasi1999}.

\subsection{Limitations}
\label{subsec:limits}

Several simplifications limit the present model. First, the 
intersection dynamics are not explicitly modelled; vehicles leave 
one segment and re-enter another without a delay. Real intersections 
(traffic lights, yield signs) introduce additional criticality 
sources. Second, the road network is generated synthetically; real 
urban networks have heterogeneous speed limits, lane counts, and 
demands. Third, the simulation is run at a fixed global density; 
in reality, density varies spatially and temporally with commute 
patterns. Future work could incorporate origin-destination demand 
matrices and empirical road network data from OpenStreetMap.

% ============================================================
\section{Conclusion}
\label{sec:conclusion}

We have demonstrated that traffic jam cascades on road networks 
exhibit the canonical signatures of self-organised criticality: 
power-law avalanche size distributions, intermittent time series, 
and spontaneous organisation to a critical state without parameter 
tuning. The power-law exponent $\tau$ is a decreasing function of 
network connectivity $\langle k \rangle$, and scale-free topologies 
produce the heaviest-tailed distributions due to cascade amplification 
at hub intersections.

These findings suggest practical implications for urban traffic 
management. Reducing the connectivity of critical hub intersections 
(e.g., through tidal flow management or congestion pricing) may 
shift the cascade distribution toward higher $\tau$, reducing the 
probability of city-wide jam events. Conversely, designing road 
networks with heterogeneous connectivity—a small number of very 
high-capacity arterials—may inadvertently increase SOC susceptibility.

The framework developed here is general and can be extended to other 
networked flow systems in chemical engineering, including pipeline 
networks, process plant utility grids, and supply chain logistics, 
wherever threshold-driven cascade dynamics arise.

% ============================================================
\section*{Acknowledgements}

The author thanks Prof.\ [Professor's Name] for the course CLL798 
and for providing the project framework. The following tools were 
used in completing this project and are acknowledged per the course 
guidelines: Python (simulation and analysis), LaTeX (typesetting), 
Claude by Anthropic (project planning assistance, code scaffolding, 
LaTeX structure), GitHub Copilot (minor code suggestions). A complete 
list of AI prompts used is provided in Appendix~\ref{app:prompts}.

% ============================================================
% References
\bibliographystyle{plainnat}
\bibliography{references}

% ── Inline bibliography if .bib file not available ──────
\begin{thebibliography}{99}

\bibitem[Bak et al.(1987)]{bak1987}
Bak, P., Tang, C., and Wiesenfeld, K. (1987).
Self-organized criticality: an explanation of $1/f$ noise.
\textit{Physical Review Letters}, 59(4):381--384.

\bibitem[Bak(1996)]{bak1996}
Bak, P. (1996).
\textit{How Nature Works: The Science of Self-Organised Criticality}.
Copernicus Press, New York.

\bibitem[Barabási and Albert(1999)]{barabasi1999}
Barabási, A.-L. and Albert, R. (1999).
Emergence of scaling in random networks.
\textit{Science}, 286(5439):509--512.

\bibitem[Clauset et al.(2009)]{clauset2009}
Clauset, A., Shalizi, C. R., and Newman, M. E. J. (2009).
Power-law distributions in empirical data.
\textit{SIAM Review}, 51(4):661--703.

\bibitem[Harris et al.(2020)]{numpy}
Harris, C. R. et al. (2020).
Array programming with NumPy.
\textit{Nature}, 585:357--362.

\bibitem[Hunter(2007)]{matplotlib}
Hunter, J. D. (2007).
Matplotlib: a 2-D graphics environment.
\textit{Computing in Science \& Engineering}, 9(3):90--95.

\bibitem[Kerner and Rehborn(1998)]{kerner1998}
Kerner, B. S. and Rehborn, H. (1998).
Experimental properties of phase transitions in traffic flow.
\textit{Physical Review Letters}, 79(20):4030--4033.

\bibitem[Nagel and Schreckenberg(1992)]{nagel1992}
Nagel, K. and Schreckenberg, M. (1992).
A cellular automaton model for freeway traffic.
\textit{Journal de Physique I}, 2(12):2221--2229.

\bibitem[Schrank et al.(2021)]{schrank2021}
Schrank, D., Eisele, B., and Lomax, T. (2021).
\textit{2021 Urban Mobility Report}.
Texas A\&M Transportation Institute.

\bibitem[Hagberg et al.(2008)]{networkx}
Hagberg, A. A., Schult, D. A., and Swart, P. J. (2008).
Exploring network structure, dynamics, and function using NetworkX.
In \textit{Proceedings of the 7th Python in Science Conference}, 
pp.\ 11--15.

\end{thebibliography}

% ============================================================
\appendix

\section{Prompt Log}
\label{app:prompts}

Per course requirements, all AI prompts used in completing this 
project are listed below.

\begin{enumerate}
  \item \textbf{Claude (Anthropic), April 2026:}
        ``Explain the deliverables and prepare a step by step first 
        principles guided plan of action [for CLL798 project].''

  \item \textbf{Claude (Anthropic), April 2026:}
        ``Give project ideas considering the 5 core deliverables 
        required from the project. Evaluate the project topic on it.''

  \item \textbf{Claude (Anthropic), April 2026:}
        ``Help me build the Traffic Jam Cascade on a Road Network 
        project --- give me the full simulation plan and LaTeX outline.''

  \item \textit{[Add further prompts here as you use them during 
        the project.]}
\end{enumerate}

\section{GitHub Repository Structure}
\label{app:github}

The repository is organised as follows:

\begin{lstlisting}[language=bash, numbers=none]
traffic-soc-cll798/
|-- README.md
|-- requirements.txt
|-- simulation/
|   |-- traffic_model.py      # NaSch model + network simulation
|   |-- run_simulation.py     # driver script (topology sweep)
|-- analysis/
|   |-- power_law_fit.py      # MLE fitting + distribution plots
|   |-- connectivity.py       # connectivity sweep + visualisations
|-- data/
|   |-- *.pkl                 # pickled simulation results
|   |-- fit_results.csv       # summary table
|-- figures/
|   |-- *.pdf                 # all paper figures
|-- report/
    |-- main.tex              # this document
    |-- main.pdf              # compiled output
    |-- references.bib
\end{lstlisting}

\end{document}
